\section{Síntesis y ampliación}

La idea de Perceiver puede condensarse en una sola oración: en lugar de dejar que un raw input enorme se atienda a sí mismo por pares en cada capa, se hace que un número reducido de learned latents lea la entrada, realice un cómputo profundo de forma interna y luego use task-specific output queries para decodificar la estructura requerida. Esta reorganización facilita que attention escale a píxeles, bytes, audio y video, y campos vectoriales densos, y además convierte el problema del diseño de arquitectura en cuestiones de latent capacity, position/modality features y query construction.

\begin{enumerate}
  \item \textbf{La motivación de la investigación es la escalabilidad del sistema}: el handcraft por modalidad no puede abarcar todos los sensores ni todos los datos científicos;
  \item \textbf{La generalidad del Transformer proviene de supuestos de conectividad débiles}: non-local attention y position-as-feature se transfieren con facilidad, pero la raw self-attention es costosa;
  \item \textbf{El latent bottleneck reescribe la complejidad}: el procesamiento de la entrada pasa de $M^2$ a $MN$, y el cómputo profundo es principalmente $N^2$;
  \item \textbf{El scaling lineal tiene condiciones}: el número de latents es fijo, el output decoder es asumible y el cuello de botella no pierde información crítica para la tarea;
  \item \textbf{Perceiver y ViT optimizan objetivos distintos}: ViT cambia el image patch prior por eficiencia; Perceiver cambia supuestos más débiles por una interfaz entre dominios;
  \item \textbf{La posición no se eliminó}: la posición Fourier/learned sigue aportando la estructura real; solo deja de codificarse de forma rígida la conectividad local;
  \item \textbf{Las output queries definen la tarea}: la posición, la modalidad y el task embedding le indican al decoder qué debe leer de los latents compartidos;
  \item \textbf{Los bytes debilitan el supuesto del tokenizer}: el costo de un vocabulary unificado es una secuencia más larga y más aprendizaje de estructura;
  \item \textbf{Optical flow pone a prueba la structured output}: una query por píxel decodifica un vector bidimensional, y la synthetic-to-real transfer muestra el potencial de los sesgos débiles;
  \item \textbf{Lo cualitativo y lo cuantitativo deben usarse juntos}: EPE, PSNR, la tasa de compresión, el rendimiento y las failure slices responden, cada uno, a preguntas distintas;
  \item \textbf{Generality y speed forman un Pareto tradeoff}: en dominios conocidos y con pocos datos, los modelos especializados como convnet/RAFT todavía pueden ser superiores;
  \item \textbf{Las direcciones pendientes están bien definidas}: small-data learning, hierarchical cross-attention, entrenamiento multimodal conjunto y una decodificación de salida más eficiente.
\end{enumerate}

\begin{importantbox}{Tarea del curso: construir un experimento de entrada/salida unificada}
Elija dos tipos de datos con diferencias estructurales marcadas, por ejemplo image patches y series de tiempo. Fije el número de latents y use un Perceiver processor compartido, pero diseñe por separado las position/modality features y las output queries. Compare: un baseline especializado, ViT/Transformer, Perceiver, distintos números de latents y la variante sin codificación posicional. Reporte la calidad, las curvas de datos de entrenamiento, el wall-clock, la memoria pico, el scaling según la longitud de entrada/salida y los failure cases.
\end{importantbox}

\begin{warningbox}{Autoevaluación final: seis conclusiones que no deben confundirse}
\textbf{Sesgo débil} no equivale a \textbf{ausencia de sesgo}; \textbf{complejidad de entrada lineal} no equivale a \textbf{costo lineal de extremo a extremo}; \textbf{input permutation robustness} no equivale a \textbf{que la estructura espacial sea inútil}; \textbf{byte-level} no equivale a \textbf{cero preprocesamiento}; \textbf{attention pattern} no equivale a \textbf{explicación causal}; \textbf{interfaz de arquitectura unificada} no equivale a \textbf{que ya se haya entrenado un único modelo conjunto}.
\end{warningbox}

\subsection{Lecturas adicionales}

Se sugiere leer en el orden «baseline Transformer—cuello de botella de entrada de Perceiver—interfaz de salida de Perceiver IO—contraste por dominio»:

{\small
\begin{itemize}[nosep,leftmargin=*]
  \item \href{https://arxiv.org/abs/1706.03762}{\textbf{Attention Is All You Need}}: QKV attention y el baseline del Transformer estándar.
  \item \href{https://arxiv.org/abs/2103.03206}{\textbf{Perceiver}}: latent bottleneck, iterative attention, ImageNet y experimentos multimodales.
  \item \href{https://arxiv.org/abs/2107.14795}{\textbf{Perceiver IO}}: output queries, bytes, optical flow y structured outputs.
  \item \href{https://arxiv.org/abs/2010.11929}{\textbf{Vision Transformer}}: image patching y contraste con el Transformer visual.
  \item \href{https://arxiv.org/abs/2005.12872}{\textbf{DETR}} y \href{https://arxiv.org/abs/2006.15055}{\textbf{Slot Attention}}: antecedentes importantes de cross-attention / learned queries en visión.
  \item \href{https://arxiv.org/abs/2006.10739}{\textbf{Fourier Features}}: fundamento teórico y experimental de la codificación de coordenadas de alta frecuencia.
  \item \href{https://arxiv.org/abs/2103.06874}{\textbf{CANINE}} y \href{https://arxiv.org/abs/2105.13626}{\textbf{ByT5}}: la línea tokenization-free / byte-level language.
  \item \href{https://arxiv.org/abs/2003.12039}{\textbf{RAFT}} y \href{https://arxiv.org/abs/2104.14544}{\textbf{AutoFlow}}: sesgos de arquitectura en optical-flow y datos de synthetic training.
\end{itemize}
}

\end{document}
